\section{Zadanie 12}
\autor{Aleksander Wieczorek}
Niech P będzie macierzą przejścia spaceru losowego na cyklu długości 2025.
Znajdź najmniejszą wartość \(t\) taką, że \(p_{x,y}(t) > 0\) dla wszystkich stanów x i y.
\subsection*{Rozwiązanie}
Zauważmy od razu, że nasz łańcuch jest nieokresowy. Z wierzchołka i krawędzią nieskierowana możemy się wrócić w 2 krokach,\footnote{zakładam, że między każdą parą sąsiednich wierzchołków da się przejść w obie strony. Gdyby tak nie było dużo ciężej przeanalizować problem, bo takie t mogłoby nie istnieć} albo przejść cały cykl w 2025 krokach, więc \[o(i)=gcd(2,2025)=1\] Cykl jest spójny, łańcuch jest nieprzywiedlny , więc każdy wierzchołek ma okres 1.\\
Stąd wynika, że istnieje takie t dla którego \(\forall_{x,y} p_{x,y}(t)>0\).\\
Dodatkowo zauważmy, że wybór punktu startowego x jest nieistotny - cykl można obrócić.\\
Jeśli przejdziemy jakąś ścieżkę w k krokach to możemy przejść nią w k+2 krokach dodając przejście krawędzią "tam i z powrotem".\\
Jedynym problemem przejścia między parą wierzchołko w t krokach jest parzystość t.\\
Dla każdego wierzchołka y \(dist(x,y)\) ma inną parzystość niż \(2025-dist(x,y)\), więc między każdą parą wierzchołków istnieje najkrótsza ścieżka o każdej parzystości.\\
Wystarczy ustalić najmniejsze t dla którego wszystkie takie ścieżki mają długość \(\le\) t.
\begin{itemize}
	\item Aby przejść z x do x w nieparzystą liczbę kroków musimy obejść cały cykl, który ma długość 2025, więc nieparzyste \(t<2025\) odpada
	\item Aby przejść z x do jego sąsiada w parzystą liczbę kroków musimy obejść prawie cały cykl, więc najkrótsza taka ścieżka ma długość 2024, więc parzyste \(t<2024\) również odpadają,
\end{itemize}
Uwzględniając oba warunki najmniejszą niewykluczoną odpowiedzią jest t=2024.
\subsection*{Odpowiedź}
\[
	\boxed{\text{2024}}
\]
